\section{Limitations and Future Work}
\label{sec:limitations}

\paragraph{User-guided pivots are not publicly evaluated prospectively.}
Internal system-verification use has produced concrete cases in which stopping or
surfacing a contradiction helped real researchers refine the operational task.
Those cases contain project-specific research details and user interactions that
cannot be disclosed here. They motivate the mechanism but do not constitute a
public prospective study of clarification quality, researcher calibration, or time
to identify the binding constraint. A publishable evaluation requires consented,
prospectively designed collaboration studies with disclosure-safe tasks.

\paragraph{Contract refinement can still fail.}
Evidence and explicit authority reduce goal drift but do not eliminate it. A
Manager or operator can approve a poorly framed tradeoff, a recorded standing
intent can omit a tacit requirement, and the report-level $K_t/X_t$ projection is
distributed across several runtime surfaces rather than one atomic transaction.
Future work should test adversarial refinements, invariant-preservation checks,
disagreement handling, and rollback of harmful contract changes.

\paragraph{Verification is only as sound as its evidence boundary.}
An executable test, formal checker, benchmark, or model Reviewer can encode the
wrong property. The runtime can record and revise a verifier, but that does not make
the verifier correct. Controlled studies should separately seed implementation
bugs, specification bugs, and mismatches between them, then measure whether the
runtime localizes the faulty side rather than merely satisfying the current check.

\paragraph{Attribution and task sequence.}
The SWE-Bench Pro experiment compares the complete \systemname runtime with Direct
Copilot. Reviewer routing is adaptive, the 22 completed Waves (24 Wave IDs with two
incomplete Waves excluded) follow one task order, and
per-Wave Direct-Copilot Token/time records are unavailable. The current results
therefore characterize whole-system behavior over this sequence. Matched
frozen-state runs, randomized task orders, and randomized review routing will
separate the contributions of persistent state, role separation, and review. The
startup-to-mature comparison is observational rather than a causal learning ablation.

\paragraph{Generality across models, domains, and evaluators.}
The seven benchmark arenas use different metrics, hardware, and execution
protocols, while the mathematical analysis follows one vertical campaign. Results
are interpreted in their native units rather than combined into a universal score.
The in-progress GLM-5.2 run on Claude Code is the first evidence that the runtime
transfers across both backbone and execution surface, but it is incomplete and has
no matched Direct baseline, so it bounds nothing yet. Repeated campaigns,
additional backbones, and external domain review will measure how the runtime
transfers across task families and research standards.

\paragraph{ACE-2 is certified for a demonstrated scope, not for silicon.}
The chip result establishes functional closure of a 24-layer, two-token
Qwen2.5-0.5B W4A8 integration and canonical mapped SKY130 synthesis with static
timing. It is not routed timing, power signoff, DRC/LVS, GDS or tapeout, silicon
validation, generation beyond two tokens, or an FPGA prototype, and the
certificate records those exclusions itself. Whether the same runtime reaches
signoff-quality physical design or fabricated silicon is untested.

\paragraph{Metric calibration and model transfer.}
The efficiency analysis labels whole role calls by mission purpose; finer segment
labels would sharpen the reasoning, action, and verification factors. Reviewer
sensitivity and false-acceptance rate require randomized routing with an external
verifier, and the counterfactual reuse value $G_L$ requires paired future tasks with
and without a state update. The runtime-to-model path is currently a training
direction: held-out trajectory splits and scaffold-removal studies will measure how
much of the long-horizon control loop can be internalized by the model.

\paragraph{Scope of the paper-production case study.}
The six paper projects come from one shared research environment, campaign-hours
overlap in calendar time, stored review snapshots are model-generated, and runtime
logging evolved across projects. Within that scope, the case study establishes
end-to-end lifecycle behavior and recovery under review across six domains and two
venue formats, not paper acceptance, scientific novelty, or superiority to human
research teams. It also does not establish a measured zero-touch rate. Recorded
costs in the public data exclude utility calls without a priced event.
